\section{Methodology}

\subsection{Framework Overview}

Figure~\ref{fig:pipeline} follows one claim from retrieval to a traceable decision. The cutoff controls which evidence may enter the graph, while dependency clustering controls how that evidence is counted. Both constraints act before classification. Calibration cannot repair a prediction formed from future or duplicated evidence.

For claim $c$ and evidence pool $E$, the retained set is
\begin{equation}
E_{\leq\tau}^{(k)}=\operatorname{TopK}_{k}\{r(c,e):e\in E,\ p(e)\leq\tau,\ o(e)\leq\tau\},
\label{eq:time-filter}
\end{equation}
where $r(c,e)$ is the fused retrieval score, $p(e)$ is the reported publication time, $o(e)$ is the earliest independent archive or crawl observation, and $k=20$. The primary set excludes uncertain timestamps. The provenance stage maps the remaining items to dependency clusters $C(E_{\leq\tau}^{(k)})$ before graph encoding.

\begin{figure*}[t]
\centering
\resizebox{\textwidth}{!}{\begin{tikzpicture}[
  font=\sffamily\scriptsize,
  stage/.style={
    draw,
    very thick,
    rounded corners=2pt,
    align=center,
    minimum height=25mm,
    inner sep=3.5mm
  },
  flow/.style={-{Stealth[length=2.4mm,width=1.8mm]}, very thick, draw=tpggray},
  note/.style={
    draw=tpggray,
    rounded corners=1.5pt,
    fill=black!3,
    align=center,
    inner sep=2mm,
    font=\sffamily\scriptsize
  }
]
  \node[stage, draw=tpgblue, fill=tpgblue!8, text width=27mm] (claim) {
    \textbf{1. Claim package}\\[1mm]
    claim text + image\\
    decision cutoff $\tau$
  };

  \node[stage, draw=tpgteal, fill=tpgteal!8, text width=38mm, right=5mm of claim] (retrieve) {
    \textbf{2. Retrieve and verify time}\\[1mm]
    entities, dates, OCR, caption\\
    BM25 + BGE-M3\\
    SigLIP2 + pHash\\
    RRF $\rightarrow$ rerank\\
    publication + observation $\leq\tau$
  };

  \node[stage, draw=tpgviolet, fill=tpgviolet!8, text width=39mm, right=5mm of retrieve] (graph) {
    \textbf{3. Model provenance}\\[1mm]
    document--image--publisher graph\\
    links + publication order\\
    dependence signals\\
    bounded cluster aggregation
  };

  \node[stage, draw=tpgorange, fill=tpgorange!9, text width=39mm, right=5mm of graph] (decision) {
    \textbf{4. Decide selectively}\\[1mm]
    heterogeneous graph encoder\\
    sufficiency $\rightarrow$ veracity\\
    temperature scaling\\
    validation-fitted abstention
  };

  \draw[flow] (claim) -- (retrieve);
  \draw[flow] (retrieve) -- (graph);
  \draw[flow] (graph) -- (decision);

  \node[note, below=3mm of retrieve, text width=38mm] (cutoff) {
    \textbf{Eligibility guard:} require a pre-cutoff observation; exclude fact-check pages
  };
  \node[note, below=3mm of graph, text width=39mm] (independence) {
    \textbf{Independence guard:} copies share one bounded source cluster
  };
  \node[note, below=3mm of decision, text width=39mm] (output) {
    \textbf{Traceable output:} label or defer, probabilities, evidence IDs, cluster IDs, cutoff
  };

  \draw[-{Stealth[length=1.8mm]}, semithick, dashed, draw=tpgteal] (retrieve) -- (cutoff);
  \draw[-{Stealth[length=1.8mm]}, semithick, dashed, draw=tpgviolet] (graph) -- (independence);
  \draw[-{Stealth[length=1.8mm]}, semithick, dashed, draw=tpgorange] (decision) -- (output);
\end{tikzpicture}
}
\caption{End-to-end \systemname{} pipeline. The temporal guard removes unavailable evidence, provenance clustering prevents copies from being treated as independent corroboration, and the final record preserves the evidence trail. Selective deferral is applied to the calibrated three-class prediction and is not a fourth task label.}
\label{fig:pipeline}
\end{figure*}

\subsection{Time-Constrained Multimodal Retrieval}

For each claim, retrieval begins with named entities, temporal expressions, OCR text, and an image caption. BM25 and BGE-M3 retrieve text candidates, while SigLIP2 embeddings and perceptual hashes retrieve images. Reciprocal-rank fusion merges lexical, dense, and visual rankings. A cross-encoder reranks the top 100 candidates. Eligibility filtering then retains at most 20 items whose publication and observation times satisfy the cutoff. All model revisions, queries, candidate lists, and filtering reasons are preserved.

\subsection{Evidence-Source Provenance Graph}

The implemented graph contains nodes for claims, images, documents, publishers, and dependency clusters. Typed edges record claim--image membership, retrieved evidence, publication, cluster membership, image containment, direct links, and publication order; reverse relations are added for message passing. Clusters are estimated from exact text hashes, token-set title similarity, perceptual image distance, and direct hyperlinks. Publication order is retained as a graph relation rather than treated as proof that the earliest item is the original source.

Let $C(E)$ denote the partition of evidence documents into dependency clusters. If cluster $j$ contains $n_j$ documents, document $i$ in that cluster receives base mass $w_i=1/n_j$. A learned within-cluster aggregator may redistribute that mass, but the cluster's total contribution is bounded. The model then combines information across clusters. Ten copies of one report therefore cannot receive the same base weight as ten independently sourced reports. Cluster quality is evaluated on a manual pair sample before the main experiment.

Source dependence is kept separate from source credibility. Domain identity can define a publisher node but does not exclude or downweight an item in the primary model. This choice preserves contradictory evidence while preventing repeated publication chains from inflating corroboration.

\subsection{Hierarchical Decision and Abstention}

A relation-aware graph encoder represents the temporal-provenance graph. The sufficiency head estimates $p(s=1\mid x,E_{\leq\tau})$ using separately audited sufficiency targets. Cases predicted as sufficient pass to a veracity head for supported or refuted; the remainder receive the insufficient-evidence label. A flat three-class model and a parameter-matched cluster-pooling model test whether the hierarchy and graph are necessary.

Let $z_s$ be the sufficiency logit and $(z_{+},z_{-})$ the supported/refuted logits. The implemented three-class distribution is
\begin{align}
p(y=\mathrm{supported}) &= \sigma(z_s)\,\operatorname{softmax}(z_{+},z_{-})_{+},\\
p(y=\mathrm{refuted}) &= \sigma(z_s)\,\operatorname{softmax}(z_{+},z_{-})_{-},\\
p(y=\mathrm{insufficient}) &= 1-\sigma(z_s).
\end{align}
Training minimizes class-weighted negative log likelihood. When a separately audited sufficiency target is available, a masked binary cross-entropy term is added with validation-selected weight; unlabeled sufficiency rows do not contribute to that auxiliary term. Evidence-node feature dropout is applied during full-model training and is removed in a named ablation.

Temperature scaling is fitted to validation logits, and a threshold selected on the same split targets 80\% answer coverage. Deferral can apply to any class when calibrated confidence is below the threshold. Each prediction records probabilities, selected evidence, dependency clusters, cutoff, timestamp policy, and whether the system answered or deferred.

